% Provenance (never printed): see sources/outline_sources.md (Year 2: tholeiitic, calc-alkaline
% and alkaline series; the TAS diagram; settings; isotopic signatures of mantle and crust).
% Classification from BGS RR 99-06 (after Streckeisen 1976, Le Bas et al. 1986, Le Maitre et
% al. 1989); series and data from Vermeesch and Pease (2021) and the GeoplotR package;
% curves and counts from figdata/bachelor-2/series.py; exercise and problem values in the
% Book 4 scratch calculators calc/ch09.py and calc/ch09pb.py.
\chapter{Séries magmáticas e contextos tectônicos}\label{ch:b2:magmatic-series}

Entregue a uma petróloga uma análise química sem mapa, e ela muitas vezes saberá dizer se a
lava foi emitida em uma dorsal, em uma ilha oceânica ou acima de uma placa em subducção. As lavas
da Islândia e as dos vulcões das Cascatas, no oeste da América do Norte, vão ambas
do basalto ao riolito; analisadas às centenas, elas seguem dois caminhos diferentes à medida
que evoluem, um cada vez mais rico em ferro, o outro não. Este capítulo começa por
nomear as rochas ígneas, por seus minerais quando eles são visíveis e por sua
química quando não são; em seguida, agrupa-as em séries, famílias de rochas
ligadas pela diferenciação, e liga cada série a um contexto tectônico, com os
diagramas que os petrólogos usam para ler um contexto a partir de uma análise, e seus limites.
% ledger: vp21:series

\begin{recall}
\textbf{Do volume do ensino médio}: basalto, andesito, riolito, gabro e granito;
rochas félsicas e máficas; texturas vulcânicas e plutônicas. \textbf{Do volume do primeiro
ano}: dorsais meso-oceânicas, zonas de subducção e arcos vulcânicos, pontos quentes e grandes
províncias ígneas, riftes continentais. \textbf{Do \cref{ch:b2:elements-in-the-earth}}:
MORB e OIB, \omterm{def:b2:elements-in-the-earth:normalised}{diagramas normalizados} ao \omterm{def:b2:elements-in-the-earth:bse}{manto primitivo} (aranhogramas), \omterm{def:b2:elements-in-the-earth:partition}{elementos compatíveis} e
incompatíveis. \textbf{Do \cref{ch:b2:isotope-geochemistry}}: o \omterm{def:b2:isotope-geochemistry:reservoirs}{manto
empobrecido} e os reservatórios isotópicos enriquecidos. \textbf{Do
\cref{ch:b2:magmatic-differentiation}}: \omterm{def:b2:magmatic-differentiation:crystallisation}{cristalização fracionada}, \omterm{def:b2:magmatic-differentiation:lld}{linha de
descida do líquido}, \omterm{def:b2:magmatic-differentiation:harker}{diagramas de Harker} e lei de Rayleigh.
\end{recall}

\section{Nomear as rochas ígneas}\label{sec:b2:magmatic-series:naming}

Uma rocha de granulação grossa mostra seus minerais e é nomeada a partir deles. Uma lava muitas vezes tem granulação
fina demais ou é vítrea demais para isso, e é nomeada a partir de sua análise química. Os dois
esquemas são convenções internacionais: eles descrevem, e nada implicam sobre o modo como
a rocha se formou.
% ledger: bgs:tas-use

\begin{definition}[Composição modal; diagrama QAPF]\label{def:b2:magmatic-series:qapf}
A \emph{composição modal}\index{composição modal} de uma rocha é a proporção em
volume de cada um de seus minerais. O \emph{diagrama QAPF}\index{diagrama QAPF} nomeia uma
rocha ígnea de granulação grossa com menos de \qty{90}{\percent} de minerais máficos a partir de
quatro componentes modais recalculados para uma soma de 100: quartzo (Q), feldspato alcalino
(A), plagioclásio (P) e os feldspatoides (F). Como quartzo e feldspatoides não
cristalizam juntos, toda rocha cai em um de dois triângulos, QAP ou FAP, unidos
ao longo de seu lado comum AP.
\end{definition}
% ledger: bgs:qapf

\begin{omfigure}
\centering
\begin{tikzpicture}[font=\scriptsize, x=1cm, y=1cm]
\def\L{6.4}
\pgfmathsetmacro{\H}{\L*0.8660254}
% upper point (q, p): q = Q fraction, p = P/(A+P)
\newcommand{\qapfup}[2]{({(#1)*\L/2 + (#2)*(1-(#1))*\L}, {(#1)*\H})}
\newcommand{\dn}[2]{({(#1)*\L/2 + (#2)*(1-(#1))*\L}, {-(#1)*\H})}
\draw[thick] (0,0) -- (\L/2,\H) -- (\L,0) -- (\L/2,-\H) -- cycle;
\draw (0,0) -- (\L,0);
\foreach \q in {0.9,0.6,0.2,0.05} \draw \qapfup{\q}{0} -- \qapfup{\q}{1};
\foreach \q in {0.1,0.6} \draw \dn{\q}{0} -- \dn{\q}{1};
\foreach \p in {0.1,0.65,0.9} \draw \qapfup{0.6}{\p} -- \qapfup{0}{\p};
\draw[dashed] \qapfup{0.6}{0.35} -- \qapfup{0.2}{0.35};
\draw \qapfup{0.2}{0.35} -- \qapfup{0}{0.35};
\foreach \p in {0.1,0.35,0.65,0.9} \draw \dn{0}{\p} -- \dn{0.1}{\p};
\foreach \p in {0.1,0.5,0.9} \draw \dn{0.1}{\p} -- \dn{0.6}{\p};
\node[anchor=south] at (\L/2,\H) {\normalsize Q};
\node[anchor=east] at (0,0) {\normalsize A};
\node[anchor=west] at (\L,0) {\normalsize P};
\node[anchor=north] at (\L/2,-\H) {\normalsize F};
\node at \qapfup{0.94}{0.5} {1a};
\node at \qapfup{0.75}{0.5} {1b};
\node[font=\tiny] at \qapfup{0.4}{0.05} {2};
\node at \qapfup{0.4}{0.225} {3a};
\node at \qapfup{0.4}{0.5} {3b};
\node at \qapfup{0.4}{0.775} {4};
\node[font=\tiny] at \qapfup{0.4}{0.95} {5};
\foreach \p/\n in {0.225/7*,0.5/8*,0.775/9*} \node at \qapfup{0.125}{\p} {\n};
\node[font=\tiny] at \qapfup{0.125}{0.05} {6*};
\node[font=\tiny] at \qapfup{0.125}{0.95} {10*};
\foreach \p/\n in {0.225/7,0.5/8,0.775/9} \node at \qapfup{0.025}{\p} {\n};
\foreach \p/\n in {0.225/7',0.5/8',0.775/9'} \node at \dn{0.05}{\p} {\n};
\node[font=\tiny] at \dn{0.3}{0.05} {11};
\node at \dn{0.35}{0.3} {12};
\node at \dn{0.35}{0.7} {13};
\node[font=\tiny] at \dn{0.3}{0.95} {14};
\node at \dn{0.75}{0.5} {15};
\node[anchor=east] at \qapfup{0.9}{0} {90};
\node[anchor=east] at \qapfup{0.6}{0} {60};
\node[anchor=east] at \qapfup{0.2}{0} {20};
\node[anchor=east] at \dn{0.1}{0} {10};
\node[anchor=east] at \dn{0.6}{0} {60};
\foreach \p/\n in {0.1/10,0.35/35,0.65/65,0.9/90} \node[anchor=south, fill=white, inner sep=0.5pt] at \qapfup{0.6}{\p} {\tiny\n};
\node[anchor=west, align=left, font=\scriptsize] at (\L+0.9,1.6) {Triângulo superior: \% Q\\em Q + A + P;\\
  retas em que\\$100\,P/(A+P)$\\vale 10, 35, 65, 90};
\node[anchor=west, align=left, font=\scriptsize] at (\L+0.9,-1.7) {Triângulo inferior: \% F\\em F + A + P;\\
  retas em que\\$100\,P/(A+P)$\\vale 10, 50, 90};
\end{tikzpicture}
\par\smallskip{\small A classificação QAPF das rochas ígneas de granulação grossa. 1a
quartzolito, 1b granitoide rico em quartzo, 2 granito de feldspato alcalino, 3a sienogranito e
3b monzogranito (juntos, granito), 4 granodiorito, 5 tonalito; 6*--10* as
variedades com quartzo de 6--10: 6 sienito de feldspato alcalino, 7 sienito, 8 monzonito, 9 monzodiorito
ou monzogabro, 10 diorito, gabro ou anortosito; 6'--10' suas variedades com
feldspatoide; 11 foide-sienito, 12 foide-monzossienito, 13 foide-monzodiorito ou
foide-monzogabro, 14 foide-diorito ou foide-gabro, 15 foidolito. Os campos estreitos 6, 6', 10 e 10'
ao longo das margens não são numerados. Segundo o esquema do British Geological Survey (baseado
em Streckeisen 1976).}
\end{omfigure}
% ledger: bgs:qapf

O campo 10 contém três nomes. O \omterm{def:b2:magmatic-series:qapf}{diagrama QAPF} não separa diorito, gabro e
anortosito, nem as rochas gabroicas entre si; outros diagramas, baseados no
plagioclásio e nos minerais máficos, fazem isso.

\begin{definition}[Diagrama álcalis totais--sílica]\label{def:b2:magmatic-series:tas}
O \emph{diagrama álcalis totais--sílica}\index{diagrama álcalis totais--sílica} (diagrama
TAS) nomeia uma rocha vulcânica cuja moda não pode ser medida, e o faz a partir de suas porcentagens
em massa de \ce{Na2O + K2O} em função de \ce{SiO2}, depois de a análise ter sido
recalculada para \qty{100}{\percent}, livre de água e de dióxido de carbono. Ele é dividido
em quinze campos por retas cujas interseções têm coordenadas fixas
(a figura depois do próximo parágrafo).
\end{definition}
% ledger: bgs:tas, bgs:tas-use

O diagrama tem três grupos de campos. Ao longo da base, com poucos álcalis, ficam os
nomes da série comum, do picrobasalto e do basalto ao riolito, separados em 45, 52,
57 e \qty{63}{\percent} de sílica e, para o riolito, por uma reta inclinada. Acima deles, os campos S e o traquito contêm rochas mais ricas em
álcalis para a mesma sílica; os campos U, o fonolito e o foidito contêm rochas tão ricas em
álcalis, e tão pobres em sílica, que cristalizam feldspatoides. Alguns campos são
subdivididos por um segundo critério: um traquibasalto é chamado de havaíto se $\ce{Na2O} - 2
\ge \ce{K2O}$, e de traquibasalto potássico no caso contrário; um tefrito com mais de
\qty{10}{\percent} de olivina em sua norma (ver adiante) é um basanito; um riolito ou um traquito cuja
razão molecular $(\ce{Na2O + K2O})/\ce{Al2O3}$ excede 1 é chamado de peralcalino. O esquema
exige análises com menos de \qty{2}{\percent} de água e \qty{0.5}{\percent} de
dióxido de carbono: a alteração acrescenta água e desloca os álcalis, e um \omterm{def:b2:magmatic-differentiation:cumulate}{cumulato} não é um
líquido.
% ledger: bgs:tas, bgs:tas-use

\begin{omfigure}
\centering
\begin{tikzpicture}
\begin{axis}[omrecord, width=0.86\linewidth, height=8.6cm,
  xmin=37, xmax=80, ymin=0, ymax=16, xtick={37,41,...,77}, ytick={0,2,...,16},
  xlabel={\ce{SiO2} (\unit{\percent})}, ylabel={\ce{Na2O + K2O} (\unit{\percent})},
  legend style={font=\scriptsize, at={(0.98,0.98)}, anchor=north east},
  tick label style={font=\scriptsize}, label style={font=\small}]
\addplot[only marks, mark=*, mark size=0.55pt, blue!70!black, opacity=0.6] table[x=SiO2, y=alk] {figdata/out/bachelor-2-series-cath-th.dat};
\addlegendentry{Islândia}
\addplot[only marks, mark=*, mark size=0.55pt, red!75!black, opacity=0.6] table[x=SiO2, y=alk] {figdata/out/bachelor-2-series-cath-ca.dat};
\addlegendentry{Cascatas}
\foreach \f in {Pc,B,O1,O2,O3,R,S1,S2,S3,T,U1,U2,U3,Ph,F}
  \addplot[black, thin, forget plot] table[x=SiO2, y=alk] {figdata/out/bachelor-2-series-tas-\f.dat};
\node[font=\tiny] at (axis cs:43,1.6) {Pc};
\node[font=\tiny, fill=white, inner sep=0.5pt] at (axis cs:50.5,0.8) {B};
\node[font=\tiny] at (axis cs:54.5,2.2) {O1};
\node[font=\tiny] at (axis cs:60,2.2) {O2};
\node[font=\tiny] at (axis cs:66.5,2.2) {O3};
\node[font=\tiny] at (axis cs:74.5,11) {R};
\node[font=\tiny, fill=white, inner sep=0.5pt] at (axis cs:49,5.8) {S1};
\node[font=\tiny] at (axis cs:52.5,7.4) {S2};
\node[font=\tiny] at (axis cs:57.5,8.9) {S3};
\node[font=\tiny] at (axis cs:64.5,12) {T};
\node[font=\tiny] at (axis cs:44,6.6) {U1};
\node[font=\tiny] at (axis cs:48.7,9.4) {U2};
\node[font=\tiny] at (axis cs:52.5,11.8) {U3};
\node[font=\tiny] at (axis cs:57,14.6) {Ph};
\node[font=\tiny] at (axis cs:39,10) {F};
\end{axis}
\end{tikzpicture}
\par\smallskip{\small O diagrama TAS (segundo Le Bas e colaboradores, 1986, tal como dado pelo
British Geological Survey), com 456 lavas e rochas intrusivas da Islândia (azul) e
580 das Cascatas (vermelho), óxidos tais como analisados. Pc picrobasalto, B basalto, O1 andesito
basáltico, O2 andesito, O3 dacito, R riolito; S1 traquibasalto, S2 traquiandesito
basáltico, S3 traquiandesito, T traquito ou traquidacito; U1 tefrito ou
basanito, U2 fonotefrito, U3 tefrifonolito, Ph fonolito, F foidito.}
\end{omfigure}
% ledger: bgs:tas, geoplotr:tas, data:b4-cath


\begin{definition}[Saturação em sílica]\label{def:b2:magmatic-series:saturation}
Um \omterm{def:b2:magma-emplacement:magma}{magma} ou uma rocha é \emph{supersaturado em sílica}\index{saturação em sílica} se tem
mais sílica do que seus outros óxidos podem combinar em minerais silicáticos, de modo que
aparece quartzo (ou outro mineral de sílica); \emph{subsaturado em sílica} se tem
pouca demais para transformar todos os seus álcalis em feldspato, de modo que aparece em seu lugar um feldspatoide (nefelina,
leucita), nunca junto com quartzo; e \emph{saturado em sílica}
entre os dois.
\end{definition}

Para uma lava cujos minerais não podem ser contados, a saturação é julgada sobre uma \emph{norma}: um
cálculo padronizado que distribui os óxidos da análise entre minerais ideais
em uma ordem fixa, os álcalis primeiro no feldspato, e indica o excesso de sílica como
quartzo, ou a falta como feldspatoide. A mais usada, a norma CIPW, separa os
dois tipos de basalto: um basalto com nefelina em sua norma é um basalto alcalino; um
sem nefelina é um basalto subalcalino, e entre estes os basaltos toleíticos têm
hiperstênio (ortopiroxênio) na norma e, muitas vezes, piroxênio pobre em cálcio na moda.
O \cref{exo:b2:magmatic-series:11} aplica uma versão simplificada dela.
% ledger: bgs:tas-use, bgs:series-terms

\section{Séries magmáticas}\label{sec:b2:magmatic-series:series}

As lavas de um vulcão, ou de uma província, raramente são iguais, mas são aparentadas:
derivam de \omterm{def:b2:magma-emplacement:magma}{magmas} parentais semelhantes pela diferenciação do
\cref{ch:b2:magmatic-differentiation}, e traçam juntas uma linha de descida.

\begin{definition}[Série magmática]\label{def:b2:magmatic-series:series}
Uma \emph{série magmática}\index{série magmática} é um conjunto de rochas ígneas de um mesmo
ciclo magmático ou de uma mesma província, ligadas pela diferenciação a partir de \omterm{def:b2:magma-emplacement:magma}{magmas} parentais semelhantes e
que partilham traços químicos e mineralógicos característicos. Três são comumente
distinguidas. A \emph{série toleítica}\index{série toleítica} parte de
basaltos saturados ou ligeiramente \omterm{def:b2:magmatic-series:saturation}{supersaturados em sílica}, com hiperstênio na norma,
e mostra um enriquecimento marcado em ferro no início de sua diferenciação. A
\emph{série cálcio-alcalina}\index{série cálcio-alcalina} vai do basalto, passando pelo
andesito, ao dacito e ao riolito, com pouco ou nenhum enriquecimento em ferro, e com
minerais hidratados (hornblenda, biotita) em seus membros evoluídos. A \emph{série
alcalina}\index{série alcalina} parte de basaltos com nefelina normativa e evolui
pelos campos S e U do diagrama TAS em direção ao traquito e ao fonolito.
\end{definition}
% ledger: bgs:series-terms, bgs:tas-use

Os nomes descrevem conjuntos, não rochas isoladas: um basalto isolado pode pertencer a qualquer
um deles, e os mesmos nomes foram usados com sentidos diferentes, de modo que uma petróloga
deve dizer o que entende por eles. A distinção entre as \omterm{def:b2:magmatic-series:series}{séries toleítica} e
cálcio-alcalina repousa sobre o ferro, e o diagrama que a mostra melhor representa três
grupos de óxidos.
% ledger: bgs:series-terms

\begin{definition}[Diagrama AFM]\label{def:b2:magmatic-series:afm}
Um \emph{diagrama AFM}\index{diagrama AFM} é um diagrama triangular cujos vértices são
$A = \ce{Na2O + K2O}$, $F = \ce{FeO_T}$ (todo o ferro como \ce{FeO}) e $M = \ce{MgO}$, em
porcentagem em massa recalculada de modo que $A + F + M = 100$. Uma linha traçada nele separa o
campo toleítico, do lado de F, do campo cálcio-alcalino.
\end{definition}
% ledger: vp21:series, geoplotr:afm

Um basalto que se diferencia perde magnésio mais depressa que qualquer outra coisa, já que a olivina e o
piroxênio são mais ricos em magnésio que o líquido, e conserva seus álcalis, que são
incompatíveis: sua linha de descida se afasta de M, em direção a A. O que acontece com o
ferro decide a série. Em um \omterm{def:b2:magma-emplacement:magma}{magma} toleítico, o ferro só sai com os
silicatos ferrosos ricos em magnésio, mais lentamente que o magnésio, de modo que a razão
$F/M$ aumenta e a tendência primeiro sobe em direção a F antes de se voltar para A: a
\emph{tendência de Fenner}. Em um \omterm{def:b2:magma-emplacement:magma}{magma} cálcio-alcalino, mais oxidado, o ferro sai também na
magnetita, que contém ferro férrico, e $F/M$ fica quase constante enquanto o líquido
vai direto em direção a A: a \emph{tendência de Bowen}. A figura abaixo mostra as duas com
rochas reais, os toleítos da Islândia e as lavas cálcio-alcalinas das Cascatas. A
linha divisória, traçada a olho em 1971, separa-as bem onde são basálticas; com
muitos álcalis, as duas tendências convergem e a linha diz pouco.
% ledger: vp21:series, geoplotr:afm, data:b4-cath

\begin{omfigure}
\centering
\begin{tikzpicture}
\begin{axis}[name=afm, width=0.5\linewidth, height=0.5\linewidth, axis equal image,
  hide axis, xmin=-0.04, xmax=1.04, ymin=-0.06, ymax=0.92, clip=false]
\addplot[only marks, mark=*, mark size=0.5pt, blue!70!black, opacity=0.6] table[x=afm_x, y=afm_y] {figdata/out/bachelor-2-series-cath-th.dat};
\addplot[only marks, mark=*, mark size=0.5pt, red!75!black, opacity=0.6] table[x=afm_x, y=afm_y] {figdata/out/bachelor-2-series-cath-ca.dat};
\addplot[very thick, black] table[x=x, y=y] {figdata/out/bachelor-2-series-afm-line.dat};
\draw[thick] (axis cs:0,0) -- (axis cs:1,0) -- (axis cs:0.5,0.8660254) -- cycle;
\node[anchor=south] at (axis cs:0.5,0.866) {F};
\node[anchor=north east] at (axis cs:0,0) {A};
\node[anchor=north west] at (axis cs:1,0) {M};
\node[font=\scriptsize, fill=white, inner sep=1pt] at (axis cs:0.52,0.62) {toleítica};
\node[font=\scriptsize, fill=white, inner sep=1pt] at (axis cs:0.33,0.12) {cálcio-alcalina};
\end{axis}
\begin{axis}[omrecord, at={(afm.east)}, anchor=west, xshift=1.3cm,
  width=0.46\linewidth, height=0.46\linewidth,
  xmin=-2.5, xmax=2.2, ymin=-4.5, ymax=1.5,
  xlabel={$\ln(A/F)$}, ylabel={$\ln(M/F)$},
  legend style={font=\scriptsize, at={(0.98,0.98)}, anchor=north east},
  tick label style={font=\scriptsize}, label style={font=\small}]
\addplot[only marks, mark=*, mark size=0.5pt, blue!70!black, opacity=0.6] table[x=lnAF, y=lnMF] {figdata/out/bachelor-2-series-lr-th.dat};
\addlegendentry{Islândia}
\addplot[only marks, mark=*, mark size=0.5pt, red!75!black, opacity=0.6] table[x=lnAF, y=lnMF] {figdata/out/bachelor-2-series-lr-ca.dat};
\addlegendentry{Cascatas}
\addplot[thick, blue!50!black, forget plot] table[x=x, y=y_th] {figdata/out/bachelor-2-series-lr-fits.dat};
\addplot[thick, red!50!black, forget plot] table[x=x, y=y_ca] {figdata/out/bachelor-2-series-lr-fits.dat};
\end{axis}
\end{tikzpicture}
\par\smallskip{\small À esquerda: o \omterm{def:b2:magmatic-series:afm}{diagrama AFM} das rochas da Islândia (azul) e das Cascatas (vermelho)
da figura TAS, com a linha divisória de Irvine e Baragar (1971). À direita:
as mesmas rochas em coordenadas log-razão, em que cada série se torna uma faixa quase
reta; as retas são regressões ortogonais, de inclinações $-1.10$ (Islândia) e $-0.72$
(Cascatas), que se encontram em $(-3.84, 2.12)$, como constataram Vermeesch e Pease (2021).}
\end{omfigure}
% ledger: geoplotr:afm, data:b4-cath, vp21:model

Das rochas islandesas dentro do alcance da linha divisória, 351 caem do lado
toleítico e 23 do outro; das rochas das Cascatas, 417 do lado cálcio-alcalino
e 57 do outro. Uma curva é uma fronteira desajeitada, e um modelo simples explica
a forma das tendências. Suponha-se que um \omterm{def:b2:magma-emplacement:magma}{magma} perca seus álcalis, seu ferro e seu magnésio
a taxas proporcionais às quantidades que contém, $\mathrm dX/\mathrm dt = -\lambda_X X$.

\begin{proposition}[Tendências em log-razão]\label{prop:b2:magmatic-series:logratio}
Se $A$, $F$ e $M$ decaem exponencialmente com constantes $\lambda_A$, $\lambda_F$ e
$\lambda_M$, o \omterm{def:b2:magma-emplacement:magma}{magma} se desloca ao longo de uma reta no plano de $\ln(A/F)$ e
$\ln(M/F)$, de inclinação $(\lambda_F - \lambda_M)/(\lambda_F - \lambda_A)$, enquanto se desloca
ao longo de uma curva no triângulo AFM.
\end{proposition}

\begin{proof}
Com $X = X_0 e^{-\lambda_X t}$, $\ln(A/F) = \ln(A_0/F_0) + (\lambda_F - \lambda_A)t$
e $\ln(M/F) = \ln(M_0/F_0) + (\lambda_F - \lambda_M)t$: ambos são lineares em $t$, logo
um é linear no outro, com a inclinação indicada. As razões $A/F$ e $M/F$ fixam o
ponto no triângulo, cujas coordenadas $A/(A + F + M)$, \dots\ não são lineares em $t$.
\end{proof}

A inclinação mais forte das rochas islandesas significa um $\lambda_F$ menor: o ferro sai
do \omterm{def:b2:magma-emplacement:magma}{magma} toleítico mais lentamente (\cref{exo:b2:magmatic-series:10}). Vermeesch e Pease mostraram que as duas
faixas da figura acima se encontram em um ponto, que eles leem como uma fonte comum
hipotética.
% ledger: vp21:model

As \omterm{def:b2:magmatic-series:series}{séries toleítica} e cálcio-alcalina diferem também em potassa. O
diagrama \ce{K2O}--\ce{SiO2} divide os basaltos, andesitos, dacitos e riolitos com
mais de \qty{48}{\percent} de sílica em tipos de baixo K, médio K e alto K por duas
retas (a figura abaixo). A maior parte das duas suítes aqui é de médio K: o
diagrama diz quanta potassa uma série carrega, não qual série ela é.
% ledger: bgs:k2o-sio2, data:b4-cath

\begin{omfigure}
\centering
\begin{tikzpicture}
\begin{axis}[omrecord, width=0.78\linewidth, height=6.2cm,
  xmin=45, xmax=79, ymin=0, ymax=5,
  xlabel={\ce{SiO2} (\unit{\percent})}, ylabel={\ce{K2O} (\unit{\percent})},
  legend style={font=\scriptsize, at={(0.02,0.98)}, anchor=north west},
  tick label style={font=\scriptsize}, label style={font=\small}]
\addplot[only marks, mark=*, mark size=0.55pt, blue!70!black, opacity=0.6] table[x=SiO2, y=K2O] {figdata/out/bachelor-2-series-cath-th.dat};
\addlegendentry{Islândia}
\addplot[only marks, mark=*, mark size=0.55pt, red!75!black, opacity=0.6] table[x=SiO2, y=K2O] {figdata/out/bachelor-2-series-cath-ca.dat};
\addlegendentry{Cascatas}
\addplot[thick, black, forget plot] table[x=x, y=y_low] {figdata/out/bachelor-2-series-k2o-lines.dat};
\addplot[thick, black, forget plot] table[x=x, y=y_high] {figdata/out/bachelor-2-series-k2o-lines.dat};
\draw[thick] (axis cs:48,0) -- (axis cs:48,0.3);
\node[font=\scriptsize, fill=white, inner sep=1pt] at (axis cs:60,3.7) {alto K};
\node[font=\scriptsize, fill=white, inner sep=1pt] at (axis cs:76.5,2.2) {médio K};
\node[font=\scriptsize, fill=white, inner sep=1pt] at (axis cs:60,0.35) {baixo K};
\end{axis}
\end{tikzpicture}
\par\smallskip{\small A potassa em função da sílica para as rochas da Islândia e das Cascatas, com
as retas por (48, 0.3)--(68, 1.2) e (48, 1.2)--(68, 3.1) que dividem os tipos de baixo K,
médio K e alto K (segundo Peccerillo e Taylor, 1976, tal como adotadas pelo
esquema internacional), prolongadas aqui além de 68.}
\end{omfigure}
% ledger: bgs:k2o-sio2, data:b4-cath

Por que a oxidação difere? Os \omterm{def:b2:magma-emplacement:magma}{magmas} cálcio-alcalinos se formam acima das zonas de subducção, onde
a placa que se desidrata fornece espécies oxidadas à cunha mantélica; os \omterm{def:b2:magma-emplacement:magma}{magmas}
cristalizam magnetita, e seus membros evoluídos hornblenda ou biotita. Os \omterm{def:b2:magma-emplacement:magma}{magmas}
toleíticos se formam nas dorsais, onde o manto funde por descompressão
(\cref{ch:b2:mantle-melting}), e são mais reduzidos.
% ledger: vp21:series

\section{Séries e contextos}\label{sec:b2:magmatic-series:settings}

Cada contexto da tectônica de placas produz seus próprios \omterm{def:b2:magma-emplacement:magma}{magmas}, porque cada um funde o
manto à sua maneira: por descompressão sob as dorsais e os pontos quentes, por adição de água
acima das zonas de subducção (\cref{ch:b2:mantle-melting}).

As \emph{dorsais meso-oceânicas} emitem basaltos toleíticos, os MORB do
\cref{ch:b2:elements-in-the-earth}, pobres em \omterm{def:b2:elements-in-the-earth:partition}{elementos incompatíveis} e com a
assinatura isotópica empobrecida de sua fonte. A Islândia, uma dorsal sobre um ponto quente, acrescenta
seus próprios toleítos a essa lista, com uma suíte completa até o riolito.
% ledger: vp21:series

As \emph{ilhas oceânicas} crescem por etapas. Um vulcão havaiano emite lavas alcalinas quando
jovem, depois constrói seu escudo de basalto toleítico, estimado em 80 a
\qty{95}{\percent} de seu volume final, em seus primeiros 1 a 2 milhões de anos; à medida que
se afasta do ponto quente, suas lavas voltam a ser alcalinas, terminando em uma longa cauda de
raras erupções alcalinas. Menos fusão dá \omterm{def:b2:magma-emplacement:magma}{magmas} mais alcalinos: os álcalis, sendo
incompatíveis, são menos diluídos em uma pequena \omterm{def:b2:mantle-melting:fraction}{fração de fundido}. Isotopicamente, os \omterm{def:b2:elements-in-the-earth:morb-oib}{basaltos de ilhas
oceânicas} amostram os reservatórios enriquecidos do \cref{ch:b2:isotope-geochemistry}.
% ledger: usgs14:hawaii-stages

Os \emph{arcos vulcânicos} emitem \omterm{def:b2:magmatic-series:series}{séries cálcio-alcalinas}. Seus \omterm{def:b2:elements-in-the-earth:trace}{elementos-traço} trazem uma
marca distintiva (a figura abaixo): normalizados ao \omterm{def:b2:elements-in-the-earth:bse}{manto primitivo}, os basaltos
de arco são enriquecidos em Rb, Ba, Th, U, Pb e Sr, elementos que a placa acrescenta à
cunha, e empobrecidos em nióbio e tântalo em relação a seus vizinhos. O basalto de arco
mediano do conjunto de dados tem $\mathrm{Nb}/\mathrm{La} = 0.49$, contra 0.97 para os MORB
e 1.23 para os \omterm{def:b2:elements-in-the-earth:morb-oib}{basaltos de ilhas oceânicas}.
% ledger: data:b4-basalt-settings

\begin{omfigure}
\centering
\begin{tikzpicture}
\begin{axis}[omrecord, width=0.86\linewidth, height=6.4cm, ymode=log,
  xmin=0.5, xmax=22.5, ymin=1, ymax=300,
  xtick={1,...,22}, xticklabels={Rb,Ba,Th,U,Nb,Ta,La,Ce,Pb,Sr,Nd,Zr,Hf,Sm,Eu,Ti,Gd,Dy,Y,Er,Yb,Lu},
  ylabel={amostra / manto primitivo},
  legend style={font=\scriptsize, at={(0.98,0.98)}, anchor=north east},
  tick label style={font=\scriptsize}, label style={font=\small}]
\addplot[thick, mark=square*, mark size=1.4pt, orange!85!black] table[x=idx, y=oib] {figdata/out/bachelor-2-series-spider.dat};
\addlegendentry{basalto de ilha oceânica}
\addplot[thick, mark=*, mark size=1.4pt, blue!70!black] table[x=idx, y=morb] {figdata/out/bachelor-2-series-spider.dat};
\addlegendentry{MORB}
\addplot[thick, mark=triangle*, mark size=1.7pt, red!75!black] table[x=idx, y=iab] {figdata/out/bachelor-2-series-spider.dat};
\addlegendentry{basalto de arco insular}
\end{axis}
\end{tikzpicture}
\par\smallskip{\small Composições medianas normalizadas ao \omterm{def:b2:elements-in-the-earth:bse}{manto primitivo} de 241 MORB, 259
\omterm{def:b2:elements-in-the-earth:morb-oib}{basaltos de ilhas oceânicas} e 256 basaltos de arcos insulares (a compilação de Vermeesch, 2006;
\omterm{def:b2:elements-in-the-earth:bse}{manto primitivo} de McDonough e Sun, 1995). Os basaltos de arco mostram a depressão em Nb
e Ta e picos em Ba, Pb e Sr; os \omterm{def:b2:elements-in-the-earth:morb-oib}{basaltos de ilhas oceânicas} são enriquecidos em todos os
elementos mais incompatíveis.}
\end{omfigure}
% ledger: data:b4-basalt-settings, ms95:pm-trace

No interior dos continentes, o \omterm{def:b2:magma-emplacement:magma}{magma} basáltico encontra uma crosta espessa que pode aquecer e fundir,
e o resultado pode ser um vulcanismo com dois modos.

\begin{definition}[Vulcanismo bimodal]\label{def:b2:magmatic-series:bimodal}
O \emph{vulcanismo bimodal}\index{vulcanismo bimodal} é a erupção, em uma mesma província,
de dois grupos de \omterm{def:b2:magma-emplacement:magma}{magmas}, basálticos e riolíticos, com poucas composições intermediárias
entre eles.
\end{definition}

O vulcanismo ao longo do traço do ponto quente de Yellowstone é bimodal: na
Planície do Rio Snake Oriental, riolitos mais antigos com uns \qty{75}{\percent} de sílica são cobertos por
basaltos mais jovens com uns \qty{50}{\percent}. A lacuna entre eles é a pista: a
diferenciação de um basalto sozinho teria deixado lavas intermediárias, e uma lacuna
aponta para uma segunda fonte para os riolitos, como a crosta que o basalto aquece.
% ledger: yvo23:bimodal

As \emph{grandes províncias ígneas}, vistas no volume do primeiro ano, emitem basalto em volumes
que nenhum outro contexto iguala; sua série é lida com os mesmos diagramas.

\section{Diagramas de discriminação}\label{sec:b2:magmatic-series:discrimination}

Os padrões da seção anterior sugerem ler um contexto diretamente a partir de uma análise,
e os petrólogos traçaram muitos diagramas para isso.

\begin{definition}[Diagrama de discriminação]\label{def:b2:magmatic-series:discrimination}
Um \emph{diagrama de discriminação}\index{diagrama de discriminação} é um diagrama de
variáveis químicas, em geral \omterm{def:b2:elements-in-the-earth:trace}{elementos-traço} imóveis, dividido em campos traçados
em torno das composições de rochas de contextos conhecidos, usado para atribuir um contexto a uma
rocha de origem desconhecida.
\end{definition}

O clássico é o triângulo Ti--Zr--Y de Pearce e Cann (1973), traçado para basaltos. Ele
usa elementos pouco deslocados pela alteração, de modo que pode ser aplicado a lavas antigas
e metamorfizadas cujo contexto se perdeu. A figura abaixo o testa com os 538 basaltos da
compilação que têm os três elementos. A maioria dos \omterm{def:b2:elements-in-the-earth:morb-oib}{basaltos de ilhas oceânicas} cai no
campo intraplaca D, 152 de 173, e a maioria dos MORB no campo B, 142 de 164; mas o campo B é
partilhado com os basaltos de arco, e 101 dos 201 basaltos de arco também caem nele.
% ledger: geoplotr:tizry, data:b4-basalt-settings

\begin{omfigure}
\centering
\begin{tikzpicture}
\begin{axis}[width=0.62\linewidth, height=0.55\linewidth, axis equal image,
  hide axis, xmin=-0.04, xmax=1.04, ymin=-0.06, ymax=0.92, clip=false,
  legend style={font=\scriptsize, at={(1.0,0.95)}, anchor=north west}]
\addplot[only marks, mark=*, mark size=0.7pt, blue!70!black, opacity=0.7] table[x=x, y=y] {figdata/out/bachelor-2-series-tizry-morb.dat};
\addlegendentry{MORB}
\addplot[only marks, mark=square*, mark size=0.7pt, orange!85!black, opacity=0.7] table[x=x, y=y] {figdata/out/bachelor-2-series-tizry-oib.dat};
\addlegendentry{basalto de ilha oceânica}
\addplot[only marks, mark=triangle*, mark size=0.9pt, red!75!black, opacity=0.7] table[x=x, y=y] {figdata/out/bachelor-2-series-tizry-iab.dat};
\addlegendentry{basalto de arco insular}
\foreach \f in {A,B,C,D}
  \addplot[thick, black, forget plot] table[x=x, y=y] {figdata/out/bachelor-2-series-tizry-\f.dat};
\draw[thick] (axis cs:0,0) -- (axis cs:1,0) -- (axis cs:0.5,0.8660254) -- cycle;
\node[anchor=south] at (axis cs:0.5,0.866) {Ti/100};
\node[anchor=north east] at (axis cs:0,0) {Zr};
\node[anchor=north west] at (axis cs:1,0) {Y $\times$ 3};
\draw[thin] (axis cs:0.33,0.36) -- (axis cs:0.22,0.46) node[anchor=south east, font=\small\bfseries, inner sep=1pt] {D};
\draw[thin] (axis cs:0.36,0.12) -- (axis cs:0.24,0.06) node[anchor=north east, font=\small\bfseries, inner sep=1pt] {C};
\draw[thin] (axis cs:0.56,0.37) -- (axis cs:0.76,0.50) node[anchor=west, font=\small\bfseries, inner sep=1pt] {A};
\draw[thin] (axis cs:0.53,0.21) -- (axis cs:0.72,0.13) node[anchor=north west, font=\small\bfseries, inner sep=1pt] {B};
\end{axis}
\end{tikzpicture}
\par\smallskip{\small O diagrama Ti--Zr--Y de Pearce e Cann (1973) com os basaltos
de contexto conhecido. A toleítos de arco insular; B MORB, toleítos de arco insular e
basaltos cálcio-alcalinos; C basaltos cálcio-alcalinos; D basaltos intraplaca. Ti a partir de
\ce{TiO2}; os três em ppm.}
\end{omfigure}
% ledger: geoplotr:tizry, data:b4-basalt-settings

\begin{method}[Usar um diagrama de discriminação]\label{met:b2:magmatic-series:discrimination}
(1) Verifique que a rocha é do tipo para o qual o diagrama foi construído, em geral um basalto, e
não um \omterm{def:b2:magmatic-differentiation:cumulate}{cumulato} ou uma lava evoluída, cujos \omterm{def:b2:elements-in-the-earth:trace}{elementos-traço} refletem a cristalização tanto
quanto a fonte. (2) Prefira elementos imóveis na alteração (Ti, Zr, Y, Nb, as terras raras
pesadas) para rochas alteradas ou metamorfizadas. (3) Represente uma suíte de amostras, não uma só:
um único ponto perto de uma fronteira não decide nada. (4) Leia o resultado como uma verossimilhança:
os campos foram traçados em torno da maioria, não da totalidade, de suas amostras de referência, e alguns campos
são partilhados. (5) Confirme com evidências independentes: outros diagramas, isótopos, o
contexto geológico.
\end{method}

A etapa (4) pode ser quantificada (\cref{exo:b2:magmatic-series:8}). Se um basalto cai no campo B, e os
três contextos fossem igualmente comuns, a probabilidade de que seja um MORB é de apenas cerca de 0.62.
Os classificadores estatísticos treinados com essas compilações herdam o mesmo limite: um basalto
guarda a memória de sua fonte e de sua fusão, e dois contextos podem produzir \omterm{def:b2:magma-emplacement:magma}{magmas}
semelhantes.
% ledger: data:b4-basalt-settings

\begin{inthelab}[Verificar uma análise de rocha total]
Antes de lançar uma rocha em qualquer desses diagramas, os petrólogos verificam a análise. Os
óxidos maiores, medidos por exemplo por fluorescência de raios X, devem somar perto de 100
com os voláteis: água e dióxido de carbono, medidos diretamente ou juntos como a
perda ao fogo, a massa que um pó perde quando fortemente aquecido. Uma grande perda ao
fogo marca uma rocha alterada, cujos álcalis podem ter se deslocado; o esquema do diagrama
TAS rejeita análises com mais de \qty{2}{\percent} de água ou
\qty{0.5}{\percent} de dióxido de carbono. As análises são então recalculadas para 100, livres de
voláteis, o ferro é expresso como \ce{FeO} total ou repartido entre ferroso e férrico,
e os materiais de referência analisados na mesma corrida são comparados com seus valores
certificados (\cref{ch:b2:field-lab-methods-2}).
\end{inthelab}
% ledger: bgs:tas-use

\section{Exercícios}

\begin{exercise}[$\star$]\label{exo:b2:magmatic-series:1}
Nomeie as seguintes rochas vulcânicas no diagrama TAS (\ce{SiO2}, \ce{Na2O}, \ce{K2O}
em \unit{\percent}, já livres de voláteis): (a) 49.60, 2.22, 0.51 (o basalto havaiano
de referência BHVO-2); (b) 58.5, 3.9, 1.6; (c) 72.0, 4.0, 4.2; (d) 47.0, 4.5, 2.5.
Para (d), dê as duas rochas possíveis e diga o que decide entre elas.
\end{exercise}
% ledger: usgs:bhvo2-major, bgs:tas

\begin{exercise}[$\star$]\label{exo:b2:magmatic-series:2}
Uma lava dá \ce{SiO2} 51.60, \ce{TiO2} 1.20, \ce{Al2O3} 16.10, \ce{FeO_T} 9.80,
\ce{MgO} 6.60, \ce{CaO} 8.20, \ce{Na2O} 3.50, \ce{K2O} 1.40, \ce{H2O} 1.40 e
\ce{CO2} \qty{0.20}{\percent}. (a) A análise é aceitável para o diagrama TAS?
(b) Recalcule-a livre de voláteis. (c) Nomeie a rocha antes e depois do
recálculo. (d) Ela é de baixo, médio ou alto K?
\end{exercise}
% ledger: bgs:tas-use, bgs:k2o-sio2

\begin{exercise}[$\star$]\label{exo:b2:magmatic-series:3}
Nomeie no \omterm{def:b2:magmatic-series:qapf}{diagrama QAPF}: (a) uma rocha de granulação grossa com \qty{25}{\percent} de quartzo,
\qty{30}{\percent} de feldspato alcalino, \qty{35}{\percent} de plagioclásio, \qty{8}{\percent} de
biotita e \qty{2}{\percent} de hornblenda; (b) uma com \qty{5}{\percent} de quartzo,
\qty{5}{\percent} de feldspato alcalino, \qty{55}{\percent} de plagioclásio e
\qty{35}{\percent} de clinopiroxênio.
\end{exercise}
% ledger: bgs:qapf

\begin{exercise}[$\star$]\label{exo:b2:magmatic-series:4}
Classifique como de baixo, médio ou alto K, com as retas da figura \ce{K2O}--\ce{SiO2}:
(a) 50.0, 0.40; (b) 56.0, 1.50; (c) 62.0, 2.90; (d) 66.0, 1.00 (\ce{SiO2},
\ce{K2O}, \unit{\percent}).
\end{exercise}
% ledger: bgs:k2o-sio2

\begin{exercise}[$\star\star$]\label{exo:b2:magmatic-series:5}
(a) Calcule $A$, $F$ e $M$ para um basalto com \ce{FeO_T} 14.5, \ce{MgO} 4.5 e
\ce{Na2O + K2O} \qty{3.2}{\percent}, e para um andesito com 6.1, 3.8 e 5.05.
(b) A linha divisória passa por $(F, A, M) = (54.5, 23.0, 22.5)$, $(52.5, 30.0,
17.5)$ e $(49.6, 35.0, 15.4)$, sendo o primeiro seu ponto mais rico em ferro; o
lado toleítico é o lado de F. De que lado cai cada rocha? (c) Que
mineral, cristalizando cedo, mantém o andesito onde ele está?
\end{exercise}
% ledger: geoplotr:afm, vp21:series

\begin{exercise}[$\star\star$]\label{exo:b2:magmatic-series:6}
O MORB mediano da compilação tem \qty{4}{ppm} de Nb e \qty{4.12}{ppm} de La; o
\omterm{def:b2:elements-in-the-earth:morb-oib}{basalto de ilha oceânica} mediano, 38 e 30.8; o basalto de arco mediano, 3 e 6.1. Com os
valores do \omterm{def:b2:elements-in-the-earth:bse}{manto primitivo} Nb 0.658 e La \qty{0.648}{ppm}, calcule os valores normalizados
e $\mathrm{Nb_N}/\mathrm{La_N}$ para cada um. Qual tem uma anomalia de nióbio, e por quê?
\end{exercise}
% ledger: data:b4-basalt-settings, ms95:pm-trace

\begin{exercise}[$\star\star$]\label{exo:b2:magmatic-series:7}
O Ti representa $0.5993$ da massa do \ce{TiO2}. Calcule as coordenadas no diagrama Ti--Zr--Y
(Ti/100, Zr, $3\mathrm Y$, reconduzidos a 100) de três basaltos: (a) \ce{TiO2}
\qty{1.6}{\percent}, Zr 110, Y \qty{32}{ppm}; (b) 2.8, 220, 29; (c) 0.8, 50, 20. Situe-os
na figura e diga o que cada posição permite concluir.
\end{exercise}
% ledger: geoplotr:tizry

\begin{exercise}[$\star\star$]\label{exo:b2:magmatic-series:8}
No teste Ti--Zr--Y deste capítulo, o campo B contém 142 de 164 MORB, 4 de 173
\omterm{def:b2:elements-in-the-earth:morb-oib}{basaltos de ilhas oceânicas} e 101 de 201 basaltos de arco. (a) Que fração de cada grupo cai
em B? (b) Se os três contextos fossem igualmente comuns, qual é a probabilidade de que um
basalto que cai em B seja um MORB? (c) O que o campo B sozinho permite dizer?
\end{exercise}
% ledger: data:b4-basalt-settings

\begin{exercise}[$\star\star$]\label{exo:b2:magmatic-series:9}
Um riolito tem \ce{Na2O} 5.5, \ce{K2O} 4.5 e \ce{Al2O3} \qty{10.5}{\percent} ($M$ =
61.98, 94.20 e \qty{101.96}{g/mol}). Calcule sua razão molecular $(\ce{Na2O +
K2O})/\ce{Al2O3}$. (a) Ele é peralcalino? (b) O feldspato alcalino contém um
óxido alcalino por \ce{Al2O3}: que fração dos álcalis não pode entrar nele?
\end{exercise}
% ledger: bgs:tas-use

\begin{exercise}[$\star\star\star$]\label{exo:b2:magmatic-series:10}
No modelo da proposição sobre as tendências em log-razão, tome $\lambda_A = 0.4$ e $\lambda_M =
0.8$. (a) Calcule a inclinação da tendência para $\lambda_F = 0.5$ e para $\lambda_F =
0.7$. (b) As rochas islandesas têm uma inclinação de $-1.10$ e as das Cascatas
$-0.72$: que $\lambda_F$ cada uma implica? (c) Que série perde seu ferro mais depressa,
e que mineral o explica? (d) O que acontece com $F/M$ ao longo de cada tendência?
\end{exercise}
% ledger: vp21:model

\begin{exercise}[$\star\star\star$]\label{exo:b2:magmatic-series:11}
Uma norma simplificada. Trabalhe em mols de óxido. Ponha todo o \ce{K2O} no ortoclásio
(\ce{K2O.Al2O3.6SiO2}), todo o \ce{Na2O} na albita (\ce{Na2O.Al2O3.6SiO2}), o
\ce{Al2O3} restante com a mesma quantidade de \ce{CaO} na anortita (\ce{CaO.Al2O3.2SiO2}), o
\ce{CaO} restante com a mesma quantidade de $(\ce{Mg,Fe})\ce{O}$ no diopsídio
(\ce{CaO.(Mg,Fe)O.2SiO2}) e o resto do $(\ce{Mg,Fe})\ce{O}$ no hiperstênio
(\ce{(Mg,Fe)O.SiO2}). Se faltar sílica, transforme o hiperstênio em olivina
(\ce{2(Mg,Fe)O.SiO2}); se ainda faltar, transforme a albita em nefelina
(\ce{Na2O.Al2O3.2SiO2}). Aplique-a a (a) BHVO-2: \ce{SiO2} 49.60, \ce{Al2O3} 13.44,
\ce{FeO_T} 11.15, \ce{MgO} 7.257, \ce{CaO} 11.40, \ce{Na2O} 2.219, \ce{K2O}
\qty{0.513}{\percent}; (b) um basalto com 45.5, 15.0, 11.5, 8.5, 10.0, 3.8 e 1.5 (massas
molares \ce{SiO2} 60.08, \ce{Al2O3} 101.96, \ce{FeO} 71.84, \ce{MgO} 40.30, \ce{CaO}
56.08, \ce{Na2O} 61.98, \ce{K2O} \qty{94.20}{g/mol}). Cada um é saturado em sílica? A
que série pertence cada basalto?
\end{exercise}
% ledger: usgs:bhvo2-major, bgs:series-terms

\begin{exercise}[$\star\star\star$]\label{exo:b2:magmatic-series:12}
Na compilação, $\mathrm{Nb}/\mathrm{La} < 0.7$ vale para 116 de 141 basaltos de arco,
21 de 101 MORB e 3 de 137 \omterm{def:b2:elements-in-the-earth:morb-oib}{basaltos de ilhas oceânicas} analisados para ambos. (a) Calcule as
três frações. (b) Com proporções iguais dos três contextos, qual é a
probabilidade de que um basalto com $\mathrm{Nb}/\mathrm{La} < 0.7$ seja um basalto de arco? (c)
Refaça o cálculo para uma região onde \qty{60}{\percent} dos basaltos são MORB, \qty{30}{\percent}
\omterm{def:b2:elements-in-the-earth:morb-oib}{basaltos de ilhas oceânicas} e \qty{10}{\percent} basaltos de arco. (d) Que lição (c)
traz para o uso dos \omterm{def:b2:magmatic-series:discrimination}{diagramas de discriminação}?
\end{exercise}
% ledger: data:b4-basalt-settings

\section{Problema: um vulcão, três lavas}

\begin{problem}[{Problema de fim de semana --- de três análises ao contexto de um vulcão e
aos cristais que ele deixou para trás}]
\label{pb:b2:magmatic-series:1}
Três lavas de um estratovulcão (inventado), de seus derrames de flanco mais antigos ao seu domo
mais jovem, dão (\unit{\percent}):

\smallskip\noindent{\small
\begin{tabular}{@{}lrrrrrrrrrr@{}}
 & \ce{SiO2} & \ce{TiO2} & \ce{Al2O3} & \ce{FeO_T} & \ce{MgO} & \ce{CaO} & \ce{Na2O} & \ce{K2O} & \ce{H2O} & \ce{CO2}\\
lava 1 & 50.4 & 0.95 & 17.6 & 8.6 & 7.9 & 10.2 & 2.9 & 0.62 & 0.70 & 0.10\\
lava 2 & 58.6 & 0.80 & 17.4 & 6.1 & 3.8 & 6.9 & 3.9 & 1.15 & 1.10 & 0.05\\
lava 3 & 65.3 & 0.55 & 16.2 & 4.0 & 1.6 & 4.3 & 4.4 & 1.85 & 1.20 & 0.05\\
\end{tabular}}

\smallskip\noindent e, em ppm: Rb 10, 20, 32; Th 1.1, 2.1, 3.4; Nb 2.0, 3.6, 5.6;
La 7, 13, 20; Sr 520, 480, 380; Zr 70, 125, 190; Y 18, 22, 24. \omterm{def:b2:elements-in-the-earth:bse}{Manto primitivo}: Nb
0.658, La 0.648, Th \qty{0.0795}{ppm}.

\noindent\textbf{Parte I --- Nomear as lavas.}
\begin{enumerate}
\item Some cada análise. As três são aceitáveis para o diagrama TAS?
\item Recalcule a lava 1 livre de voláteis: dê o fator, \ce{SiO2} e
\ce{Na2O + K2O}.
\item Faça o mesmo para as lavas 2 e 3.
\item Nomeie as três lavas no diagrama TAS.
\item Os nomes mudariam sem o recálculo?
\item Uma rocha de granulação grossa com a composição da lava 2 cristalizou nas raízes do
vulcão. Que esquema a nomeia, e em que campo você a esperaria?
\end{enumerate}

\noindent\textbf{Parte II --- Que série?}
\begin{enumerate}[resume]
\item Classifique cada lava como de baixo, médio ou alto K.
\item Calcule $A$, $F$, $M$ para cada uma (livres de voláteis).
\item A linha divisória passa por $(F, A, M) = (42.8, 17.1, 40.1)$, $(46.9, 18.1,
35.0)$, $(52.5, 30.0, 17.5)$, $(49.6, 35.0, 15.4)$, $(40.8, 50.0, 9.2)$ e $(37.9,
55.0, 7.1)$, com o campo toleítico do lado de F. Interpolando $F$ linearmente em
$A$ entre pontos vizinhos, de que lado cai cada lava?
\item Calcule \ce{FeO_T}/\ce{MgO} para cada lava. Há enriquecimento em ferro?
\item A que série pertence o vulcão? Que mineral e que condições
o explicam?
\end{enumerate}

\noindent\textbf{Parte III --- Que contexto?}
\begin{enumerate}[resume]
\item Normalize Nb, La e Th da lava 1 ao \omterm{def:b2:elements-in-the-earth:bse}{manto primitivo}.
\item Calcule $\mathrm{Nb_N}/\mathrm{La_N}$ para a lava 1 e compare-o com as medianas
do \cref{exo:b2:magmatic-series:6}.
\item Calcule $\mathrm{Th_N}/\mathrm{Nb_N}$. O que trouxe o tório?
\item Calcule as coordenadas Ti--Zr--Y da lava 1 (\ce{TiO2} da análise livre de
voláteis). O que diz seu campo? Por que não lançar as lavas 2 e 3?
\item Conclua: onde cresceu esse vulcão?
\end{enumerate}

\noindent\textbf{Parte IV --- Quanto cristalizou?}
\begin{enumerate}[resume]
\item A potassa não entra em nenhum dos primeiros minerais. A partir do \ce{K2O} livre de voláteis das
lavas 1 e 2, calcule a fração $F$ de líquido restante quando a lava 2 se formou, se ela
deriva da lava 1 por \omterm{def:b2:magmatic-differentiation:crystallisation}{cristalização fracionada}.
\item Faça o mesmo com o rubídio e com o tório.
\item O zircônio tem um coeficiente global $D \approx 0.1$. Use a lei de Rayleigh para estimar $F$ a partir
dele.
\item O estrôncio cai da lava 2 para a lava 3. Calcule $F$ a partir da potassa entre as lavas 2 e
3, e use-o para estimar o $D$ global do Sr nesse intervalo. Que mineral é
responsável?
\item Da lava 1 à lava 3, que fração do líquido cristalizou?
\item Se o \omterm{def:b2:magma-emplacement:magma}{magma} tivesse assimilado crosta granítica rica em potassa, a estimativa pela
potassa seria alta demais ou baixa demais?
\item Faça a média das quatro estimativas da Parte IV para as lavas 1 a 2 e indique a massa de
cristais deixada para trás por \qty{100}{kg} de basalto para fazer o andesito, o número-chave
deste problema.
\end{enumerate}
\end{problem}
% ledger: bgs:tas-use, bgs:tas, bgs:k2o-sio2, geoplotr:afm, geoplotr:tizry, ms95:pm-trace
